\section{三个群元素直接给出六边长及体积}

这是人类 review 要求的计算路线：从三个群元素出发，经六条边长直接得到体积。
设 $q(p)$ 为群元的四元数坐标，因
$q(p)\cdot q(r)=\tfrac12\Tr(p^{-1}r)$，六条短测地线边长按用户边序 $(01,02,03,23,13,12)$ 为
\begin{equation}
\boxed{\begin{aligned}
\ell_1=\ell_{01}&=\arccos\tfrac12\Tr g_1,&
\ell_2=\ell_{02}&=\arccos\tfrac12\Tr(g_1g_2),\\
\ell_3=\ell_{03}&=\arccos\tfrac12\Tr(g_1g_2g_3),&
\ell_4=\ell_{23}&=\arccos\tfrac12\Tr g_3,\\
\ell_5=\ell_{13}&=\arccos\tfrac12\Tr(g_2g_3),&
\ell_6=\ell_{12}&=\arccos\tfrac12\Tr g_2 .
\end{aligned}}\label{eq:group-six-lengths}
\end{equation}
对边为 $j$ 与 $j+3$。这些是单位 $S^3$ 上的弧长，不能代入欧氏 Cayley--Menger 体积。

这六条边长的完整体积函数，是第 \ref{sec:symbolic-ade} 节的式 \eqref{eq:exact-master-phase}。
它是 Murakami 定理 1.2\cite{murakami} 的群元代入形式，固定 $z$ 的全部偏导都已展开，
取向由 $D=\tfrac14\Tr[A(BC-CB)]$（$A=g_1$，$B=g_1g_2$，$C=g_1g_2g_3$）的符号给出。
当三个群元属于 ADE 有限群时，第 \ref{sec:polar-dual} 节的式 \eqref{eq:polar-master}
进一步给出不含 $\Li_2$ 的精确值。
六条边长不能区分镜像，$D=0$ 时也不能一概置零；完整分支见第 \ref{sec:E-explicit} 节。

\paragraph{与历史接口的边序换算。}
历史数值接口 \texttt{edge\_lengths\_quaternions} 与本报告后半部保留的原始推导使用旧边序
$(23,13,03,01,02,12)$，它依次对应上面的 $(\ell_4,\ell_5,\ell_3,\ell_1,\ell_2,\ell_6)$。
全文新结果一律使用用户边序；旧边序只出现在历史接口与原推导中。
